\documentclass[11pt]{article}

\usepackage[margin=0.85in]{geometry}
\usepackage[T1]{fontenc}
\usepackage[utf8]{inputenc}
\usepackage[strings]{underscore}
\usepackage{booktabs}
\usepackage{longtable}
\usepackage{array}
\usepackage{graphicx}
\usepackage{hyperref}
\usepackage{enumitem}
\usepackage{xcolor}
\usepackage{listings}
\usepackage[capitalise]{cleveref}

\hypersetup{colorlinks=true,linkcolor=blue,citecolor=blue,urlcolor=blue}

\lstdefinestyle{pyblock}{
  basicstyle=\ttfamily\small,
  breaklines=true,
  columns=fullflexible,
  keywordstyle=\color{blue!60!black},
  commentstyle=\color{gray},
  stringstyle=\color{green!40!black},
  frame=single,
  language=Python,
}
\lstset{style=pyblock}

\newcolumntype{L}[1]{>{\raggedright\arraybackslash}p{#1}}

\title{\textbf{aci\_lib User Guide}}
\author{ACI-CO processed-data toolkit}
\date{\today}

\begin{document}
\sloppy
\maketitle

\begin{abstract}
This is a practical how-to for the contents of this repository:
\texttt{src/aci\_lib/} (a small Python package for reading already-computed
ACI-CO pipeline outputs into pandas/xarray) and
\texttt{notebooks/aci\_lib\_quickstart.ipynb} (a Google Colab notebook built
on top of it). For \emph{why} the underlying data looks the way it does --
the ACI-CO methodology, the pipeline that produced it, which scripts wrote
which files -- see the companion \texttt{ARCHITECTURE.pdf} in this same
repository instead. This document only covers using what is already built.
\end{abstract}

\tableofcontents

%======================================================================
\section{What's in this repository}
%======================================================================

\begin{longtable}{@{}L{3.6cm}L{10.4cm}@{}}
\toprule
\textbf{Path} & \textbf{Contents} \\
\midrule
\endhead
\texttt{src/aci\_lib/} & The library itself: four small modules (\Cref{sec:api}). No third-party code, only pandas/xarray at call time. \\
\texttt{notebooks/aci\_lib\_\allowbreak quickstart.ipynb} & A ready-to-run Colab notebook: clones this repo, points \texttt{aci\_lib} at your data, browses the catalog, loads two example datasets, plots one (\Cref{sec:notebook}). \\
\texttt{ARCHITECTURE.tex} / \texttt{.pdf} & The source pipeline's architecture: scientific target, data flow, known issues. Reference material, not a how-to. \\
\bottomrule
\end{longtable}

\textbf{What is not in this repository:} the data itself. \texttt{data/processed/}
from the source ACI-CO project is roughly 13\,GB across 8{,}000+ files --
too large for git -- so it is never committed here. Every workflow below
starts by pointing \texttt{aci\_lib} at a copy of that directory that you
supply (\Cref{sec:data-supply}).

%======================================================================
\section{Setup}
%======================================================================

\subsection{Requirements}

\texttt{aci\_lib} itself has no hard dependencies beyond the standard
library at import time; individual functions import what they need lazily:

\begin{itemize}[leftmargin=*]
\item \texttt{pandas} -- always (used by \texttt{list\_datasets()}, and by \texttt{load()} for CSV/XLSX).
\item \texttt{openpyxl} -- only if you load an \texttt{.xlsx} dataset.
\item \texttt{xarray} and \texttt{dask} -- only if you load a NetCDF dataset.
\item A NetCDF backend for xarray (\texttt{netCDF4} or \texttt{h5netcdf}) -- same condition.
\item \texttt{matplotlib} -- only if you plot; not imported by the library at all.
\end{itemize}

In Google Colab, all of these are preinstalled. Locally, install them with
your usual environment manager if they are not already present (they are,
in the source ACI-CO project's environment, since its own scripts use the
same stack).

\subsection{Getting the code onto \texttt{sys.path}}

\texttt{aci\_lib} is not published as a package -- there is no
\texttt{pip install}. You put \texttt{src/} on \texttt{sys.path} and import
it directly.

\paragraph{In Colab:}
\begin{lstlisting}
!git clone --depth 1 https://github.com/mdgordillob/aca_aci_collab.git
import sys
sys.path.insert(0, "/content/aca_aci_collab/src")
import aci_lib
\end{lstlisting}

\paragraph{Locally, from a checkout of this repo:}
\begin{lstlisting}
import sys, os
sys.path.insert(0, os.path.abspath("src"))  # or "../src" from notebooks/
import aci_lib
\end{lstlisting}

%======================================================================
\section{Supplying the data}
\label{sec:data-supply}
%======================================================================

\texttt{aci\_lib} needs to know where a copy of \texttt{data/processed/}
lives. It resolves this in the following order the first time it is
needed:

\begin{enumerate}[leftmargin=*]
\item An explicit \texttt{aci\_lib.set\_data\_dir(path)} call.
\item The \texttt{ACI\_DATA\_DIR} environment variable.
\item Auto-detection: walking up from \texttt{src/aci\_lib/} to the nearest
  directory containing a \texttt{.git} folder (the repo root), then
  checking whether \texttt{data/processed/} exists under it.
\end{enumerate}

If none of these resolve to a real directory, every \texttt{aci\_lib} call
that needs data raises \texttt{FileNotFoundError} with a message pointing
back to this section.

\begin{longtable}{@{}L{3.0cm}L{5.3cm}L{5.7cm}@{}}
\toprule
\textbf{Method} & \textbf{When to use it} & \textbf{How} \\
\midrule
\endhead
Local checkout of the \emph{source} project & You have the ACI-CO pipeline repo cloned locally, with \texttt{data/processed/} already populated on disk & Nothing to do -- auto-detection (step 3 above) finds it. \\
Mount Google Drive & Colab, and you have already uploaded \texttt{data/processed/} to your Drive & \texttt{drive.mount(\allowbreak"/content/drive")}, then \texttt{aci\_lib.set\_data\_dir(\allowbreak"/content/drive/MyDrive/\allowbreak <path>/processed")} \\
Upload + unzip & Colab, one-off, no Drive setup wanted & \texttt{files.upload()} a zip, \texttt{!unzip} it, then \texttt{set\_data\_dir()} to the unzipped path \\
Environment variable & Any script/notebook where you'd rather configure via environment than code & \texttt{export ACI\_DATA\_DIR=/path/to/processed} (shell) or \texttt{os.environ["ACI\_DATA\_DIR"] = ...} before importing \\
\bottomrule
\end{longtable}

All four are demonstrated in \texttt{notebooks/aci\_lib\_quickstart.ipynb}.

%======================================================================
\section{API reference}
\label{sec:api}
%======================================================================

Four functions, all imported directly off the \texttt{aci\_lib} package.

\subsection{\texttt{set\_data\_dir(path)}}

Points the library at a \texttt{data/processed}-like directory for the
rest of the session. Raises \texttt{FileNotFoundError} immediately if
\texttt{path} is not a directory (fails fast rather than deferring the
error to the first \texttt{load()} call).

\begin{lstlisting}
aci_lib.set_data_dir("/content/drive/MyDrive/aca_data/processed")
\end{lstlisting}

\subsection{\texttt{get\_data\_dir()}}

Returns whatever directory \texttt{aci\_lib} is currently configured to
read from (see the resolution order in \Cref{sec:data-supply}). Useful to
sanity-check setup before calling \texttt{list\_datasets()}:

\begin{lstlisting}
print(aci_lib.get_data_dir())
\end{lstlisting}

\subsection{\texttt{list\_datasets(base\_dir=None, kind=None)}}

Returns a \texttt{pandas.DataFrame}, one row per loadable dataset, columns
\texttt{name}, \texttt{kind}, \texttt{n\_files}, \texttt{description}.
\texttt{base\_dir} overrides the configured data directory for this call
only (does not call \texttt{set\_data\_dir}). \texttt{kind} filters to one
of \texttt{"tabular"}, \texttt{"grid"}, \texttt{"grid\_series"},
\texttt{"image"}, \texttt{"raw"} (\Cref{sec:naming} explains what each
means).

\begin{lstlisting}
aci_lib.list_datasets()                       # everything
aci_lib.list_datasets(kind="tabular")         # only CSV/XLSX datasets
aci_lib.list_datasets()[lambda d: d.name.str.contains("colombia")]
\end{lstlisting}

\subsection{\texttt{describe(name, base\_dir=None)}}

Returns the underlying \texttt{Dataset} record for one catalog entry: its
\texttt{name}, \texttt{kind}, the full list of file paths it is built
from, and its \texttt{description}. Useful to see exactly which files a
\texttt{grid\_series} dataset will open before loading potentially
gigabytes of NetCDF.

\begin{lstlisting}
ds = aci_lib.describe("anomalias_colombia/anomalies_temperature")
print(ds.kind, ds.n_files)
print(ds.files[:3])
\end{lstlisting}

\subsection{\texttt{load(name, base\_dir=None)}}

The main entry point. Returns different Python types depending on the
dataset's \texttt{kind}:

\begin{longtable}{@{}L{2.4cm}L{4.6cm}L{7.0cm}@{}}
\toprule
\textbf{Kind} & \textbf{Returns} & \textbf{Notes} \\
\midrule
\endhead
\texttt{tabular} & \texttt{pandas.DataFrame} & \texttt{pd.read\_csv} or \texttt{pd.read\_excel} depending on the source file's extension. \\
\texttt{grid} & \texttt{xarray.Dataset} & A single NetCDF file, opened with \texttt{xr.open\_dataset}. \\
\texttt{grid\_series} & \texttt{xarray.Dataset} & Many NetCDF files (e.g.\ one per year-month) opened lazily with \texttt{xr.open\_mfdataset} and concatenated along a \texttt{file} dimension, in chronological order. Backed by dask -- see \Cref{sec:perf}. \\
\texttt{image} & \texttt{str} & The file path (a \texttt{.png}); open it with \texttt{PIL.Image.open(...)} or \texttt{plt.imread(...)} yourself. \\
\texttt{raw} & \texttt{str} & The file path for anything else (\texttt{.txt}, \texttt{.ascii}, \texttt{.zip}, ...); read it yourself. \\
\bottomrule
\end{longtable}

Raises \texttt{KeyError} if \texttt{name} is not in the catalog -- the
message points back to \texttt{list\_datasets()}.

\begin{lstlisting}
df = aci_lib.load("anomalias_colombia/anomalies_temperature_combined")
grid = aci_lib.load("anomalias_colombia/anomalies_temperature")
img_path = aci_lib.load("confusion_matrix")
\end{lstlisting}

%======================================================================
\section{How dataset names are built}
\label{sec:naming}
%======================================================================

Names mirror \texttt{data/processed/}'s directory structure, with the file
extension stripped:

\begin{itemize}[leftmargin=*]
\item A file directly under \texttt{data/processed/}, e.g.
  \texttt{ungrd\_monthly.csv}, becomes \texttt{"ungrd\_monthly"}
  (\texttt{kind="tabular"}).
\item A file in a subdirectory, e.g.
  \texttt{anomalias\_colombia/anomalies\_temperature\_combined.csv}, becomes
  \texttt{"anomalias\_colombia/anomalies\_temperature\_combined"}.
\item A \emph{run} of NetCDF files that share a prefix and differ only by a
  trailing \texttt{\_YYYY} or \texttt{\_YYYY\_MM}, e.g.\
  \texttt{anomalies\_temperature\_1961\_1.nc} through
  \texttt{\_2024\_12.nc} (768 files), collapses to one entry named after
  the shared prefix -- \texttt{"anomalias\_colombia/\allowbreak anomalies\_temperature"} --
  with \texttt{kind="grid\_series"}.
\item A single file that happens to match that same \texttt{\_YYYY[\_MM]}
  pattern but has no siblings (e.g.\ a baseline file literally named
  \texttt{aci\_temp\_1961\_1990.nc}) is \emph{not} treated as a series --
  it keeps its full name and \texttt{kind="grid"}.
\end{itemize}

When in doubt, do not guess a name -- call \texttt{aci\_lib.list\_datasets()}
and filter/search the \texttt{name} column.

%======================================================================
\section{Common recipes}
%======================================================================

\subsection{Plot a national time series}

\begin{lstlisting}
import matplotlib.pyplot as plt

df = aci_lib.load("anomalias_colombia/anomalies_temperature_combined")
df = df.sort_values(["year", "month"])
plt.plot(range(len(df)), df["t_90"])
plt.title("National T90 anomaly, monthly")
plt.show()
\end{lstlisting}

\subsection{Pull one month out of a gridded series}

\begin{lstlisting}
grid = aci_lib.load("anomalias_colombia/anomalies_temperature")
first_month = grid.isel(file=0)   # already chronologically first
first_month["t_90"].plot()
\end{lstlisting}

\subsection{Compare the same variable across regions}

\begin{lstlisting}
regions = ["antioquia", "cali", "colombia", "cundinamarca_bogota", "valle_cauca"]
series = {
    r: aci_lib.load(f"anomalias_{r}/anomalies_temperature_combined")
    for r in regions
}
\end{lstlisting}

\subsection{Only browse what is small/tabular first}

Gridded datasets can be gigabytes; when exploring interactively, start
with \texttt{kind="tabular"} to avoid triggering large lazy loads by
accident:

\begin{lstlisting}
aci_lib.list_datasets(kind="tabular")
\end{lstlisting}

%======================================================================
\section{Performance notes}
\label{sec:perf}
%======================================================================

\begin{itemize}[leftmargin=*]
\item \texttt{grid\_series} datasets are opened with dask chunking (one
  chunk per source file) and are \emph{lazy}: \texttt{aci\_lib.load(...)}
  itself does not read the underlying data into memory. Values are only
  computed when you call \texttt{.values}, \texttt{.compute()}, or plot.
\item Subset (\texttt{.isel}/\texttt{.sel}) before triggering computation,
  not after -- e.g.\ \texttt{grid.isel(file=0)["t\_90"].plot()} reads one
  file's worth of data; \texttt{grid["t\_90"].values} reads all of them.
\item \texttt{list\_datasets()} and \texttt{describe()} only walk the
  directory tree and match filenames -- they never open a data file, so
  they are cheap to call repeatedly.
\end{itemize}

%======================================================================
\section{Troubleshooting}
%======================================================================

\begin{longtable}{@{}L{5.4cm}L{8.6cm}@{}}
\toprule
\textbf{Symptom} & \textbf{Fix} \\
\midrule
\endhead
\texttt{FileNotFoundError: Could not find data/processed/} & No data source configured yet -- follow \Cref{sec:data-supply}. \\
\texttt{KeyError: No dataset named ...} & The name does not match the catalog exactly (case-sensitive, no leading/trailing \texttt{.csv}). Run \texttt{aci\_lib.list\_datasets()} and search/filter its \texttt{name} column instead of guessing. \\
\texttt{ModuleNotFoundError: No module named 'xarray'} (or \texttt{dask}, \texttt{netCDF4}) & Only triggered when loading a NetCDF dataset. Install the missing package -- not needed at all if you only load tabular datasets. \\
\texttt{load()} on a \texttt{grid\_series} is slow or uses a lot of memory & You likely triggered a full compute on an unsubsetted dataset -- see \Cref{sec:perf}. \\
Colab session restarted / variables gone & \texttt{sys.path}, \texttt{aci\_lib}'s configured data directory, and any Drive mount are session state -- re-run the setup cells at the top of the notebook. \\
\bottomrule
\end{longtable}

%======================================================================
\section{Running the Colab notebook end to end}
\label{sec:notebook}
%======================================================================

\texttt{notebooks/aci\_lib\_quickstart.ipynb} is organized to be run
top to bottom:

\begin{enumerate}[leftmargin=*]
\item \textbf{Colab detection} -- a small cell sets \texttt{IN\_COLAB} so
  the rest of the notebook behaves correctly whether it is running in
  Colab or from a local checkout.
\item \textbf{Get the code} -- clones this repo (Colab) or adds the
  sibling \texttt{../src} to \texttt{sys.path} (local).
\item \textbf{Supply the data} -- Drive-mount cell is active by default;
  the upload/unzip alternative is present but commented out. Edit the
  \texttt{DATA\_DIR} path to match where you put your data.
\item \textbf{Browse the catalog} -- calls \texttt{list\_datasets()} and
  shows the first 20 rows.
\item \textbf{Load examples} -- one tabular dataset, one gridded series,
  and a plot of the first month in that series.
\item \textbf{Reference} -- pointers to \texttt{describe()} and to
  \texttt{ARCHITECTURE.pdf} for pipeline-level context.
\end{enumerate}

Open it directly in Colab via the badge at the top of the notebook, or at:
\url{https://colab.research.google.com/github/mdgordillob/aca_aci_collab/blob/main/notebooks/aci_lib_quickstart.ipynb}

\end{document}
